\documentclass[10pt,a4paper]{amsart}
\usepackage{amsmath,amssymb,amsthm,graphicx}
% For dealing with references we use the comment environment
\usepackage{verbatim}
\newenvironment{reference}{\comment}{\endcomment}
%\newenvironment{reference}{}{}
\newenvironment{slogan}{\comment}{\endcomment}
\newenvironment{history}{\comment}{\endcomment}
% For commutative diagrams we use Xy-pic
\usepackage[all]{xy}

% We use 2cell for 2-commutative diagrams.
\xyoption{2cell}
\UseAllTwocells
\theoremstyle{plain}
\newtheorem{theorem}[subsection]{Theorem}
\newtheorem{proposition}[subsection]{Proposition}
\newtheorem{lemma}[subsection]{Lemma}

\theoremstyle{definition}
\newtheorem{definition}[subsection]{Definition}
\newtheorem{example}[subsection]{Example}
\newtheorem{exercise}[subsection]{Exercise}
\newtheorem{situation}[subsection]{Situation}

\theoremstyle{remark}
\newtheorem{remark}[subsection]{Remark}
\newtheorem{remarks}[subsection]{Remarks}

\numberwithin{equation}{subsection}
% Macros
%
\def\lim{\mathop{\mathrm{lim}}\nolimits}
\def\colim{\mathop{\mathrm{colim}}\nolimits}
\def\Spec{\mathop{\mathrm{Spec}}}
\def\Hom{\mathop{\mathrm{Hom}}\nolimits}
\def\Ext{\mathop{\mathrm{Ext}}\nolimits}
\def\SheafHom{\mathop{\mathcal{H}\!\mathit{om}}\nolimits}
\def\SheafExt{\mathop{\mathcal{E}\!\mathit{xt}}\nolimits}
\def\Sch{\mathit{Sch}}
\def\Mor{\mathop{\mathrm{Mor}}\nolimits}
\def\Ob{\mathop{\mathrm{Ob}}\nolimits}
\def\Sh{\mathop{\mathit{Sh}}\nolimits}
\def\NL{\mathop{N\!L}\nolimits}
\def\CH{\mathop{\mathrm{CH}}\nolimits}
\def\proetale{{pro\text{-}\acute{e}tale}}
\def\etale{{\acute{e}tale}}
\def\QCoh{\mathit{QCoh}}
\def\Ker{\mathop{\mathrm{Ker}}}
\def\Im{\mathop{\mathrm{Im}}}
\def\Coker{\mathop{\mathrm{Coker}}}
\def\Coim{\mathop{\mathrm{Coim}}}

% Boxtimes
%
\DeclareMathSymbol{\boxtimes}{\mathbin}{AMSa}{"02}

%
% Macros for moduli stacks/spaces
%
\def\QCohstack{\mathcal{QC}\!\mathit{oh}}
\def\Cohstack{\mathcal{C}\!\mathit{oh}}
\def\Spacesstack{\mathcal{S}\!\mathit{paces}}
\def\Quotfunctor{\mathrm{Quot}}
\def\Hilbfunctor{\mathrm{Hilb}}
\def\Curvesstack{\mathcal{C}\!\mathit{urves}}
\def\Polarizedstack{\mathcal{P}\!\mathit{olarized}}
\def\Complexesstack{\mathcal{C}\!\mathit{omplexes}}
% \Pic is the operator that assigns to X its picard group, usage \Pic(X)
% \Picardstack_{X/B} denotes the Picard stack of X over B
% \Picardfunctor_{X/B} denotes the Picard functor of X over B
\def\Pic{\mathop{\mathrm{Pic}}\nolimits}
\def\Picardstack{\mathcal{P}\!\mathit{ic}}
\def\Picardfunctor{\mathrm{Pic}}
\def\Deformationcategory{\mathcal{D}\!\mathit{ef}}
\usepackage{cleveref}
\crefname{originalequation}{original equation}{original equations}
\Crefname{originalequation}{Original equation}{Original equations}
\makeatletter
\newlabel{equation-descent-etale}{{0BTI}{source}{}{sourcepointer-1}{}}
\newlabel{equation-descent-etale@cref}{{[originalequation][0][]0BTI}{[1][0][]source}{}{}{}}
\makeatother
\crefname{originallemma}{original lemma}{original lemmas}
\Crefname{originallemma}{Original lemma}{Original lemmas}
\makeatletter
\newlabel{lemma-fully-faithful-cases}{{0BTL}{source}{}{sourcepointer-2}{}}
\newlabel{lemma-fully-faithful-cases@cref}{{[originallemma][0][]0BTL}{[1][0][]source}{}{}{}}
\makeatother
\crefname{originallemma}{original lemma}{original lemmas}
\Crefname{originallemma}{Original lemma}{Original lemmas}
\makeatletter
\newlabel{lemma-reduce-to-affine}{{0BTM}{source}{}{sourcepointer-3}{}}
\newlabel{lemma-reduce-to-affine@cref}{{[originallemma][0][]0BTM}{[1][0][]source}{}{}{}}
\makeatother
\crefname{originallemma}{original lemma}{original lemmas}
\Crefname{originallemma}{Original lemma}{Original lemmas}
\makeatletter
\newlabel{algebra-lemma-limit-integral}{{0BTG}{source}{}{sourcepointer-4}{}}
\newlabel{algebra-lemma-limit-integral@cref}{{[originallemma][0][]0BTG}{[1][0][]source}{}{}{}}
\makeatother
\crefname{originallemma}{original lemma}{original lemmas}
\Crefname{originallemma}{Original lemma}{Original lemmas}
\makeatletter
\newlabel{limits-lemma-descend-finite-presentation}{{01ZM}{source}{}{sourcepointer-5}{}}
\newlabel{limits-lemma-descend-finite-presentation@cref}{{[originallemma][0][]01ZM}{[1][0][]source}{}{}{}}
\makeatother
\crefname{originallemma}{original lemma}{original lemmas}
\Crefname{originallemma}{Original lemma}{Original lemmas}
\makeatletter
\newlabel{limits-lemma-descend-separated-finite-presentation}{{01ZQ}{source}{}{sourcepointer-6}{}}
\newlabel{limits-lemma-descend-separated-finite-presentation@cref}{{[originallemma][0][]01ZQ}{[1][0][]source}{}{}{}}
\makeatother
\crefname{originallemma}{original lemma}{original lemmas}
\Crefname{originallemma}{Original lemma}{Original lemmas}
\makeatletter
\newlabel{limits-lemma-descend-etale}{{07RP}{source}{}{sourcepointer-7}{}}
\newlabel{limits-lemma-descend-etale@cref}{{[originallemma][0][]07RP}{[1][0][]source}{}{}{}}
\makeatother
\crefname{originalproposition}{original proposition}{original propositions}
\Crefname{originalproposition}{Original proposition}{Original propositions}
\makeatletter
\newlabel{morphisms-proposition-generic-flatness}{{052A}{source}{}{sourcepointer-8}{}}
\newlabel{morphisms-proposition-generic-flatness@cref}{{[originalproposition][0][]052A}{[1][0][]source}{}{}{}}
\makeatother
\crefname{originalreference}{original reference}{original references}
\Crefname{originalreference}{Original reference}{Original references}
\makeatletter
\newlabel{more-morphisms-lemma-separated-locally-quasi-finite-morphisms-fppf-descend}{{02W8}{source}{}{sourcepointer-9}{}}
\newlabel{more-morphisms-lemma-separated-locally-quasi-finite-morphisms-fppf-descend@cref}{{[originalreference][0][]02W8}{[1][0][]source}{}{}{}}
\makeatother
\crefname{originallemma}{original lemma}{original lemmas}
\Crefname{originallemma}{Original lemma}{Original lemmas}
\makeatletter
\newlabel{descent-lemma-descending-property-quasi-compact}{{02KQ}{source}{}{sourcepointer-10}{}}
\newlabel{descent-lemma-descending-property-quasi-compact@cref}{{[originallemma][0][]02KQ}{[1][0][]source}{}{}{}}
\makeatother
\crefname{originallemma}{original lemma}{original lemmas}
\Crefname{originallemma}{Original lemma}{Original lemmas}
\makeatletter
\newlabel{descent-lemma-descending-property-separated}{{02KU}{source}{}{sourcepointer-11}{}}
\newlabel{descent-lemma-descending-property-separated@cref}{{[originallemma][0][]02KU}{[1][0][]source}{}{}{}}
\makeatother
\crefname{originallemma}{original lemma}{original lemmas}
\Crefname{originallemma}{Original lemma}{Original lemmas}
\makeatletter
\newlabel{morphisms-lemma-generically-finite}{{02NW}{source}{}{sourcepointer-12}{}}
\newlabel{morphisms-lemma-generically-finite@cref}{{[originallemma][0][]02NW}{[1][0][]source}{}{}{}}
\makeatother
\crefname{originalreference}{original reference}{original references}
\Crefname{originalreference}{Original reference}{Original references}
\makeatletter
\newlabel{more-algebra-lemma-henselization-flat}{{0AGU}{source}{}{sourcepointer-13}{}}
\newlabel{more-algebra-lemma-henselization-flat@cref}{{[originalreference][0][]0AGU}{[1][0][]source}{}{}{}}
\makeatother
\crefname{originalreference}{original reference}{original references}
\Crefname{originalreference}{Original reference}{Original references}
\makeatletter
\newlabel{more-algebra-lemma-finite-etale-equivalence}{{09ZL}{source}{}{sourcepointer-14}{}}
\newlabel{more-algebra-lemma-finite-etale-equivalence@cref}{{[originalreference][0][]09ZL}{[1][0][]source}{}{}{}}
\makeatother
\crefname{originalreference}{original reference}{original references}
\Crefname{originalreference}{Original reference}{Original references}
\makeatletter
\newlabel{more-algebra-lemma-integral-over-henselian-pair}{{09XK}{source}{}{sourcepointer-15}{}}
\newlabel{more-algebra-lemma-integral-over-henselian-pair@cref}{{[originalreference][0][]09XK}{[1][0][]source}{}{}{}}
\makeatother
\crefname{originallemma}{original lemma}{original lemmas}
\Crefname{originallemma}{Original lemma}{Original lemmas}
\makeatletter
\newlabel{lemma-split-henselian}{{0BTN}{source}{}{sourcepointer-16}{}}
\newlabel{lemma-split-henselian@cref}{{[originallemma][0][]0BTN}{[1][0][]source}{}{}{}}
\makeatother
\crefname{originalsection}{original section}{original sections}
\Crefname{originalsection}{Original section}{Original sections}
\makeatletter
\newlabel{schemes-section-glueing-schemes}{{01JA}{source}{}{sourcepointer-17}{}}
\newlabel{schemes-section-glueing-schemes@cref}{{[originalsection][0][]01JA}{[1][0][]source}{}{}{}}
\makeatother
\title{S001-proposition-effective}
\author{The Stacks Project Authors}
\date{stacks-project:a04446e57ec1fbc252a871afcec7752fb2807b14; GNU Free Documentation License 1.2 (as published in the Stacks Project COPYING)}
\begin{document}
\maketitle
\noindent\textit{Editorial note.} All mathematical statement and proof text is reproduced verbatim from the author's chapter etale.tex of the Stacks Project, fixed at commit a04446e57ec1fbc252a871afcec7752fb2807b14. The pinned preamble is supplied as its own file; the document wrapper class, the xr-hyper wiring and the \textbackslash{}externaldocument list are omitted because the assembler supplies the class.\par\medskip
\noindent\textit{Editorial note.} Numbering inside this standalone excerpt is editorial, not the full Stacks book print numbering.\par\medskip
\noindent\textit{Editorial note.} The candidate was previously judged H2 by the initial-corpus stage; that judgement and its reason are reused verbatim and were not re-derived here. Original candidate key: S001-proposition-effective.\par\medskip
\noindent\textit{Editorial note.} Reference labels in the proof resolve to navigation pointers to the original author source; these pointers do not reproduce proofs of the prerequisites.\par\medskip
\noindent\textit{Allowed standard prerequisites (editorial).}\begin{itemize}
\item Prerequisites named by the author inside the counted span. The counted span is reproduced verbatim, including the author's own citations. Its prerequisites are the earlier results of the same fixed Stacks Project version that the author cites there and in the chapter's lead-in.
\end{itemize}\medskip
\ifdefined\hypertarget\hypertarget{sourcepointer-1}{}\fi
\noindent\textit{Source reference.} Original equation 0BTI: sources/etale.tex, source lines 2035--2035. Navigation pointer used in the counted span; the span contains the pointer only and the cited result is not reproduced inside it. Registered as a non-load-bearing navigation stub in the established form for Stacks material. No bibliography citation is involved.\par\medskip
\ifdefined\hypertarget\hypertarget{sourcepointer-2}{}\fi
\noindent\textit{Source reference.} Original lemma 0BTL: sources/etale.tex, source lines 2124--2124. Navigation pointer used in the counted span; the span contains the pointer only and the cited result is not reproduced inside it. Registered as a non-load-bearing navigation stub in the established form for Stacks material. No bibliography citation is involved.\par\medskip
\ifdefined\hypertarget\hypertarget{sourcepointer-3}{}\fi
\noindent\textit{Source reference.} Original lemma 0BTM: sources/etale.tex, source lines 2158--2158. Navigation pointer used in the counted span; the span contains the pointer only and the cited result is not reproduced inside it. Registered as a non-load-bearing navigation stub in the established form for Stacks material. No bibliography citation is involved.\par\medskip
\ifdefined\hypertarget\hypertarget{sourcepointer-4}{}\fi
\noindent\textit{Source reference.} Original lemma 0BTG: sources/algebra.tex, source lines 32506--32506. Navigation pointer used in the counted span; the span contains the pointer only and the cited result is not reproduced inside it. Registered as a non-load-bearing navigation stub in the established form for Stacks material. No bibliography citation is involved.\par\medskip
\ifdefined\hypertarget\hypertarget{sourcepointer-5}{}\fi
\noindent\textit{Source reference.} Original lemma 01ZM: sources/limits.tex, source lines 2370--2370. Navigation pointer used in the counted span; the span contains the pointer only and the cited result is not reproduced inside it. Registered as a non-load-bearing navigation stub in the established form for Stacks material. No bibliography citation is involved.\par\medskip
\ifdefined\hypertarget\hypertarget{sourcepointer-6}{}\fi
\noindent\textit{Source reference.} Original lemma 01ZQ: sources/limits.tex, source lines 1851--1851. Navigation pointer used in the counted span; the span contains the pointer only and the cited result is not reproduced inside it. Registered as a non-load-bearing navigation stub in the established form for Stacks material. No bibliography citation is involved.\par\medskip
\ifdefined\hypertarget\hypertarget{sourcepointer-7}{}\fi
\noindent\textit{Source reference.} Original lemma 07RP: sources/limits.tex, source lines 1949--1949. Navigation pointer used in the counted span; the span contains the pointer only and the cited result is not reproduced inside it. Registered as a non-load-bearing navigation stub in the established form for Stacks material. No bibliography citation is involved.\par\medskip
\ifdefined\hypertarget\hypertarget{sourcepointer-8}{}\fi
\noindent\textit{Source reference.} Original proposition 052A: sources/morphisms.tex, source lines 5172--5172. Navigation pointer used in the counted span; the span contains the pointer only and the cited result is not reproduced inside it. Registered as a non-load-bearing navigation stub in the established form for Stacks material. No bibliography citation is involved.\par\medskip
\ifdefined\hypertarget\hypertarget{sourcepointer-9}{}\fi
\noindent\textit{Source reference.} Original reference 02W8: sources/more-morphisms.tex, source lines 16279--16279. Navigation pointer used in the counted span; the span contains the pointer only and the cited result is not reproduced inside it. Registered as a non-load-bearing navigation stub in the established form for Stacks material. No bibliography citation is involved.\par\medskip
\ifdefined\hypertarget\hypertarget{sourcepointer-10}{}\fi
\noindent\textit{Source reference.} Original lemma 02KQ: sources/descent.tex, source lines 5471--5471. Navigation pointer used in the counted span; the span contains the pointer only and the cited result is not reproduced inside it. Registered as a non-load-bearing navigation stub in the established form for Stacks material. No bibliography citation is involved.\par\medskip
\ifdefined\hypertarget\hypertarget{sourcepointer-11}{}\fi
\noindent\textit{Source reference.} Original lemma 02KU: sources/descent.tex, source lines 5581--5581. Navigation pointer used in the counted span; the span contains the pointer only and the cited result is not reproduced inside it. Registered as a non-load-bearing navigation stub in the established form for Stacks material. No bibliography citation is involved.\par\medskip
\ifdefined\hypertarget\hypertarget{sourcepointer-12}{}\fi
\noindent\textit{Source reference.} Original lemma 02NW: sources/morphisms.tex, source lines 13192--13192. Navigation pointer used in the counted span; the span contains the pointer only and the cited result is not reproduced inside it. Registered as a non-load-bearing navigation stub in the established form for Stacks material. No bibliography citation is involved.\par\medskip
\ifdefined\hypertarget\hypertarget{sourcepointer-13}{}\fi
\noindent\textit{Source reference.} Original reference 0AGU: sources/more-algebra.tex, source lines 3087--3087. Navigation pointer used in the counted span; the span contains the pointer only and the cited result is not reproduced inside it. Registered as a non-load-bearing navigation stub in the established form for Stacks material. No bibliography citation is involved.\par\medskip
\ifdefined\hypertarget\hypertarget{sourcepointer-14}{}\fi
\noindent\textit{Source reference.} Original reference 09ZL: sources/more-algebra.tex, source lines 3383--3383. Navigation pointer used in the counted span; the span contains the pointer only and the cited result is not reproduced inside it. Registered as a non-load-bearing navigation stub in the established form for Stacks material. No bibliography citation is involved.\par\medskip
\ifdefined\hypertarget\hypertarget{sourcepointer-15}{}\fi
\noindent\textit{Source reference.} Original reference 09XK: sources/more-algebra.tex, source lines 2753--2753. Navigation pointer used in the counted span; the span contains the pointer only and the cited result is not reproduced inside it. Registered as a non-load-bearing navigation stub in the established form for Stacks material. No bibliography citation is involved.\par\medskip
\ifdefined\hypertarget\hypertarget{sourcepointer-16}{}\fi
\noindent\textit{Source reference.} Original lemma 0BTN: sources/etale.tex, source lines 2200--2200. Navigation pointer used in the counted span; the span contains the pointer only and the cited result is not reproduced inside it. Registered as a non-load-bearing navigation stub in the established form for Stacks material. No bibliography citation is involved.\par\medskip
\ifdefined\hypertarget\hypertarget{sourcepointer-17}{}\fi
\noindent\textit{Source reference.} Original section 01JA: sources/schemes.tex, source lines 2510--2510. Navigation pointer used in the counted span; the span contains the pointer only and the cited result is not reproduced inside it. Registered as a non-load-bearing navigation stub in the established form for Stacks material. No bibliography citation is involved.\par\medskip
\begin{proposition}
\label{proposition-effective}
Let $f : X \to S$ be a surjective integral morphism.
The functor (\ref{equation-descent-etale}) induces an equivalence
$$
\begin{matrix}
\text{schemes quasi-compact,}\\
\text{separated, \'etale over }S
\end{matrix}
\longrightarrow
\begin{matrix}
\text{descent data }(V, \varphi)\text{ relative to }X/S\text{ with}\\
V\text{ quasi-compact, separated, \'etale over }X
\end{matrix}
$$
\end{proposition}

\begin{proof}
By Lemma \ref{lemma-fully-faithful-cases} the
functor (\ref{equation-descent-etale})
is fully faithful and the same remains the case after any
base change $S \to S'$. Let $(V, \varphi)$ be a descent data
relative to $X/S$ with $V \to X$ quasi-compact, separated, and \'etale.
We can use Lemma \ref{lemma-reduce-to-affine}
to see that it suffices to prove the effectivity 
Zariski locally on $S$. In particular we may and do
assume that $S$ is affine.

\medskip\noindent
If $S$ is affine we can find a directed set $\Lambda$ and
an inverse system $X_\lambda \to S_\lambda$
of finite morphisms of affine schemes of finite type over
$\Spec(\mathbf{Z})$ such that $(X \to S) = \lim (X_\lambda \to S_\lambda)$.
See Algebra, Lemma \ref{algebra-lemma-limit-integral}.
Since limits commute with limits we deduce that
$X \times_S X = \lim X_\lambda \times_{S_\lambda} X_\lambda$
and 
$X \times_S X \times_S X = \lim
X_\lambda \times_{S_\lambda} X_\lambda \times_{S_\lambda} X_\lambda$.
Observe that $V \to X$ is a morphism of finite presentation.
Using Limits, Lemmas \ref{limits-lemma-descend-finite-presentation}
we can find an $\lambda$ and a descent datum $(V_\lambda, \varphi_\lambda)$
relative to $X_\lambda/S_\lambda$ whose pullback to $X/S$ is
$(V, \varphi)$. Of course it is enough to show that
$(V_\lambda, \varphi_\lambda)$ is effective. Note that $V_\lambda$
is quasi-compact by construction.
After possibly increasing $\lambda$ we may assume
that $V_\lambda \to X_\lambda$ is separated and \'etale, see
Limits, Lemma \ref{limits-lemma-descend-separated-finite-presentation} and
\ref{limits-lemma-descend-etale}.
Thus we may assume that $f$ is finite surjective and
$S$ affine of finite type over $\mathbf{Z}$.

\medskip\noindent
Consider an open $S' \subset S$ such that the pullback $(V', \varphi')$
of $(V, \varphi)$ to $X' = X \times_S S'$ is effective. Below we will
prove, that $S' \not = S$ implies there is a strictly larger open over
which the descent datum is effective. Since $S$ is Noetherian (and hence
has a Noetherian underlying topological space) this will finish the proof.
Let $\xi \in S$ be a generic point of an irreducible component of the
closed subset $Z = S \setminus S'$.
If $\xi \in S'' \subset S$ is an open over which the descent datum is
effective, then the descent datum is effective over
$S' \cup S''$ by the glueing argument of the first paragraph. Thus
in the rest of the proof we may replace $S$ by an affine open
neighbourhood of $\xi$.

\medskip\noindent
After a first such replacement we may assume that $Z$ is irreducible
with generic point $Z$. Let us endow $Z$ with the reduced induced
closed subscheme structure. After another shrinking we may assume
$X_Z = X \times_S Z = f^{-1}(Z) \to Z$ is flat, see
Morphisms, Proposition \ref{morphisms-proposition-generic-flatness}.
Let $(V_Z, \varphi_Z)$ be the pullback of the descent datum to $X_Z/Z$.
By More on Morphisms, Lemma
\ref{more-morphisms-lemma-separated-locally-quasi-finite-morphisms-fppf-descend}
this descent datum is effective and we obtain an \'etale morphism
$U_Z \to Z$ whose base change is isomorphic to $V_Z$ in a manner
compatible with descent data.
Of course $U_Z \to Z$ is quasi-compact and separated
(Descent, Lemmas \ref{descent-lemma-descending-property-quasi-compact} and
\ref{descent-lemma-descending-property-separated}).
Thus after shrinking once more we may assume
that $U_Z \to Z$ is finite, see
Morphisms, Lemma \ref{morphisms-lemma-generically-finite}.

\medskip\noindent
Let $S = \Spec(A)$ and let $I \subset A$ be the prime ideal corresponding
to $Z \subset S$. Let $(A^h, IA^h)$ be the henselization of the pair
$(A, I)$. Denote $S^h = \Spec(A^h)$ and $Z^h = V(IA^h) \cong Z$.
We claim that it suffices to show effectivity after base change to
$S^h$. Namely, $\{S^h \to S, S' \to S\}$ is an fpqc covering
($A \to A^h$ is flat by More on Algebra, Lemma
\ref{more-algebra-lemma-henselization-flat}) and
by More on Morphisms, Lemma
\ref{more-morphisms-lemma-separated-locally-quasi-finite-morphisms-fppf-descend}
we have fpqc descent for separated \'etale morphisms.
Namely, if $U^h \to S^h$ and $U' \to S'$ are the objects
corresponding to the pullbacks $(V^h, \varphi^h)$ and
$(V', \varphi')$, then the required isomorphisms
$$
U^h \times_S S^h \to S^h \times_S V^h
\quad\text{and}\quad
U^h \times_S S' \to S^h \times_S U'
$$
are obtained by the fully faithfulness pointed out in the first
paragraph. In this way we reduce to the situation described in
the next paragraph.

\medskip\noindent
Here $S = \Spec(A)$, $Z = V(I)$, $S' = S \setminus Z$ where
$(A, I)$ is a henselian pair, we have $U' \to S'$ corresponding
to the descent datum $(V', \varphi')$ and we have a finite \'etale
morphism $U_Z \to Z$ corresponding to the descent datum
$(V_Z, \varphi_Z)$. We no longer have that $A$ is of finite type
over $\mathbf{Z}$; but the rest of the argument will not even use
that $A$ is Noetherian.
By More on Algebra, Lemma \ref{more-algebra-lemma-finite-etale-equivalence}
we can find a finite \'etale morphism $U_{fin} \to S$ whose
restriction to $Z$ is isomorphic to $U_Z \to Z$.
Write $X = \Spec(B)$ and $Y = V(IB)$. Since $(B, IB)$ is a henselian pair
(More on Algebra, Lemma \ref{more-algebra-lemma-integral-over-henselian-pair})
and since the restriction $V \to X$ to $Y$
is finite (as base change of $U_Z \to Z$) we see that
there is a canonical disjoint union decomposition
$$
V = V_{fin} \amalg V_{away}
$$
were $V_{fin} \to X$ is finite and where $V_{away}$ has no
points lying over $Y$. See Lemma \ref{lemma-split-henselian}.
Using the uniqueness of this decomposition over $X \times_S X$
we see that $\varphi$ preserves it and we obtain
$$
(V, \varphi) = (V_{fin}, \varphi_{fin}) \amalg (V_{away}, \varphi_{away})
$$
in the category of descent data.
By More on Algebra, Lemma \ref{more-algebra-lemma-finite-etale-equivalence}
there is a unique isomorphism
$$
X \times_S U_{fin} \longrightarrow V_{fin}
$$
compatible with the given isomorphism $Y \times_Z U_Z \to V \times_X Y$
over $Y$.
By the uniqueness we see that this isomorphism is compatible
with descent data, i.e.,
$(X \times_S U_{fin}, can) \cong (V_{fin}, \varphi_{fin})$.
Denote $U'_{fin} = U_{fin} \times_S S'$. By fully faithfulness
we obtain a morphism $U'_{fin} \to U'$ which is
the inclusion of an open (and closed) subscheme.
Then we set $U = U_{fin} \amalg_{U'_{fin}} U'$ (glueing of schemes as
in Schemes, Section \ref{schemes-section-glueing-schemes}).
The morphisms $X \times_S U_{fin} \to V$ and
$X \times_S U' \to V$ glue to a morphism $X \times_S U \to V$
which is the desired isomorphism.
\end{proof}

\end{document}
